\section{Transformation of matrices}

\SourcePage{54}{57}
The matrix elements of one and the same physical quantity may be defined with
respect to different sets of wave functions. These may be, for example, wave
functions of stationary states described by different sets of physical
quantities, or wave functions of stationary states of the same system placed
in different external fields. The question therefore arises of transforming
matrices from one representation to another.

Let $\psi_n(q)$ and $\psi'_n(q)$ ($n=1,2,\ldots$) be two complete systems of
orthonormal functions. They are related by a linear transformation
\begin{equation}
  \psi'_n=\sum_m S_{mn}\psi_m,
  \tag{12.1}
\end{equation}
which is simply the expansion of the functions $\psi'_n$ in the complete
system $\psi_m$. This transformation can be written in operator form as
\begin{equation}
  \psi'_n=\hat S\psi_n.
  \tag{12.2}
\end{equation}
The operator $\hat S$ must satisfy a certain condition if the functions
$\psi'_n$ are to be orthonormal whenever the functions $\psi_n$ are.
Indeed, substituting (12.2) into
$\int\psi_m'^*\psi'_n\,dq=\delta_{mn}$ and using the definition (3.14) of the
transposed operator, we obtain
\[
  \int(\hat S\psi_m)^*\hat S\psi_n\,dq
  =\int\psi_m^*\tilde{S^*}\hat S\psi_n\,dq=\delta_{mn}.
\]
\SourcePage{55}{58}
For this equality to hold for every $m,n$, we must have
$\tilde{S^*}\hat S=1$, or
\begin{equation}
  \tilde{S^*}\equiv\hat S^+=\hat S^{-1},
  \tag{12.3}
\end{equation}
i.e. the inverse operator coincides with the adjoint. Operators having this
property are called \textit{unitary}. By virtue of this property the
transformation $\psi_n=\hat S^{-1}\psi'_n$, inverse to (12.1), is given by
\begin{equation}
  \psi_n=\sum_m S_{nm}^*\psi'_m.
  \tag{12.4}
\end{equation}
Writing $\hat S^+\hat S=1$ or $\hat S\hat S^+=1$ in matrix form gives the
unitarity conditions
\begin{equation}
  \sum_l S_{lm}^*S_{ln}=\delta_{mn}
  \tag{12.5}
\end{equation}
or
\begin{equation}
  \sum_l S_{ml}^*S_{nl}=\delta_{mn}.
  \tag{12.6}
\end{equation}

Now consider a physical quantity $f$ and write its matrix elements in the
``new'' representation, i.e. with respect to the functions $\psi'_n$. They are
given by
\[
 \int\psi_m'^*\hat f\psi'_n\,dq
 =\int(\tilde{S^*}\psi_m^*)(\hat f\hat S\psi_n)\,dq
 =\int\psi_m^*\hat S^{-1}\hat f\hat S\psi_n\,dq.
\]
It follows that the matrix of $\hat f$ in the new representation coincides
with the matrix of the operator
\begin{equation}
  \hat f'=\hat S^{-1}\hat f\hat S
  \tag{12.7}
\end{equation}
in the old representation\footnote{If
$\{\hat f,\hat g\}=-i\hbar\hat c$ is the commutation rule for two operators,
then after transformation (12.7) we obtain
$\{\hat f',\hat g'\}=-i\hbar\hat c'$, so the rule is unchanged. It was noted
in the note on p.~46 that $c$ is the quantum analogue of the classical Poisson
bracket $[f,g]$. Poisson brackets in classical mechanics, however, are
invariant under canonical variable transformations (see Vol.~I, \S~45).
One may therefore say that unitary transformations in quantum mechanics
serve as the counterpart of
canonical transformations in classical mechanics.}.
\SourcePage{56}{59}
The \textit{trace} of a matrix is defined as the sum of its diagonal elements
and is written $\Tr f$\footnote{The source's notation comes from the German
\textit{Spur}, trace. The notation $\Tr$, from the English
\textit{trace}, is also used. Consideration of a matrix trace of course assumes
convergence of the sum over $n$.}:
\begin{equation}
  \Tr f=\sum_n f_{nn}.
  \tag{12.8}
\end{equation}
First note that the trace of a product of two matrices does not depend on the
order of the factors:
\begin{equation}
  \Tr(fg)=\Tr(gf).
  \tag{12.9}
\end{equation}
Indeed, by the rule for matrix multiplication,
\[
 \Tr(fg)=\sum_n\sum_k f_{nk}g_{kn}
 =\sum_k\sum_n g_{kn}f_{nk}=\Tr(gf).
\]
Similarly, the trace of a product of several matrices is unchanged by a cyclic
permutation of the factors; thus
\begin{equation}
 \Tr(fgh)=\Tr(hfg)=\Tr(ghf).
 \tag{12.10}
\end{equation}
The most important property of the trace is its independence of the system of
functions relative to which the matrix elements are defined. Indeed,
\begin{equation}
 \Tr f'=\Tr(S^{-1}fS)=\Tr(SS^{-1}f)=\Tr f.
 \tag{12.11}
\end{equation}
Also, a unitary transformation leaves invariant the sum of the squared moduli
of the transformed functions. From (12.6),
\begin{equation}
 \sum_i|\psi'_i|^2
 =\sum_{k,l,i}S_{ki}S_{li}^*\psi_k\psi_l^*
 =\sum_{k,l}\delta_{kl}\psi_k\psi_l^*=\sum_k|\psi_k|^2.
 \tag{12.12}
\end{equation}

Every unitary operator can be represented as
\begin{equation}
 \hat S=e^{i\hat R},
 \tag{12.13}
\end{equation}
where $\hat R$ is Hermitian; indeed, $\hat R^+=\hat R$ gives
$\hat S^+=e^{-i\hat R^+}=e^{-i\hat R}=\hat S^{-1}$.

We also note the expansion
\begin{equation}
 \hat f'=\hat S^{-1}\hat f\hat S
 =\hat f+\{\hat f,i\hat R\}
 +\frac12\{\{\hat f,i\hat R\},i\hat R\}+\ldots,
 \tag{12.14}
\end{equation}
\SourcePage{57}{60}
which is readily verified directly by expanding the factors
$\exp(\pm i\hat R)$ in powers of $\hat R$. This expansion can be useful when
$\hat R$ is proportional to a small parameter, so that (12.14) becomes an
expansion in powers of that parameter.
